%==============================================================================
% CHAPTER 25: Pharmaceutical Oscillatory Response
%==============================================================================

\chapter{Pharmaceutical Oscillatory Response}
\label{chap:pharmacology}

\epigraph{%
  A drug that only does one thing\\
  does not understand biology.%
}{attributed to Paul Ehrlich, \textit{Collected Papers} (1913),\\
  amended for the partition framework}

\begin{chapterbox}
Drugs are not molecular keys fitting biochemical locks; they are partition
topology modifiers that change the partition depth $\Depth_{\mathrm{pathway}}$
of specific biochemical cascades. Drug efficacy is proportional to
$|\Delta\Depth_D|$ at the target partition; drug toxicity arises from
$|\Delta\Depth_D|$ at non-target partitions; drug specificity is the ratio of
target to total $\Depth$ modification. The O$_2$--H$^+$ electromagnetic coupling
framework explains genome-wide drug effects through O$_2$ aggregation constants
($K_{\mathrm{agg}} > 10^4\,\mathrm{M}^{-1}$) that alter the O$_2$--H$^+$
resonance ($\omega_{\mathrm{O}_2}/\omega_{\mathrm{H}^+} = 1/4$) and thereby
modify DNA charge state globally, changing gene expression without sequence-
specific targeting. DNA palindromes are partition fixed points in the cardinal
map; palindrome density in pharmacogenes (CYP2D6, CYP2C19, etc.) predicts
pharmacogenomic impact per variant. The Genome Parser tool maps the
pharmacogenomic landscape using partition depth analysis, predicting
drug-response profiles from sequence alone. Three case studies (warfarin,
clopidogrel, metformin) are analysed in full partition-depth detail.
\end{chapterbox}

%------------------------------------------------------------------------------
\section{The Partition Topology View of Drug Action}
\label{sec:drug_topology}
%------------------------------------------------------------------------------

Classical pharmacology is built on the receptor-occupancy hypothesis: a drug
molecule binds to a receptor protein with affinity $K_d$, and the biological
effect is proportional to receptor occupancy. This framework predicts dose-response
curves, competitive inhibition, and allosteric effects. It does not explain:
\begin{enumerate}[label=(\roman*)]
  \item Why drugs with identical receptor binding affinities have different
        clinical efficacy profiles;
  \item Why the same drug produces dramatically different effects in different
        genetic backgrounds (pharmacogenomics);
  \item Why drugs targeting specific receptors affect gene expression across
        dozens of pathways not connected to the receptor;
  \item Why drug combinations often produce non-additive (synergistic or
        antagonistic) effects that are unpredictable from individual drug profiles.
\end{enumerate}

The partition framework addresses all four failures by treating drug-target
interaction as a partition topology modification rather than a lock-and-key
binding event.

\begin{definition}[Drug as Partition Topology Modifier]
\label{def:drug_modifier}
A drug $D$ interacting with target $T$ modifies the partition depth of the
downstream pathway $\mathrm{P}_T$ by:
\begin{equation}
  \Depth_{\mathrm{P}_T}^{(D)} = \Depth_{\mathrm{P}_T} - \Delta\Depth_D,
  \label{eq:drug_depth}
\end{equation}
where $\Delta\Depth_D > 0$ for inhibitors (reducing pathway depth, making
fewer downstream states accessible) and $\Delta\Depth_D < 0$ for activators
(increasing pathway depth). Drug efficacy, specificity, and toxicity are
defined by:
\begin{align}
  \text{Efficacy:}   &\quad \mathcal{E}_D \propto |\Delta\Depth_D(\text{target})| \label{eq:efficacy} \\
  \text{Toxicity:}   &\quad \mathcal{T}_D \propto \sum_{P \neq P_T} |\Delta\Depth_D(P)| \label{eq:toxicity} \\
  \text{Specificity:}&\quad \mathcal{S}_D = \frac{|\Delta\Depth_D(P_T)|}{\sum_P |\Delta\Depth_D(P)|}. \label{eq:specificity}
\end{align}
\end{definition}

\begin{remark}[Relation to Classical Pharmacology]
In the classical framework, efficacy is $E_{\max} \times [D] / (K_d + [D])$
(Hill equation), and selectivity is the ratio of $K_d$ values at target vs.\
off-target receptors. These are consistent with Equations~\eqref{eq:efficacy}--\eqref{eq:specificity}
when the partition depth modification is proportional to receptor occupancy:
$\Delta\Depth_D = \Delta\Depth_{\max} \times [D]/(K_d + [D])$.
The partition framework adds the term $\sum_{P \neq P_T}|\Delta\Depth_D(P)|$
(total off-target depth modification) which is not accessible from $K_d$
measurements alone: two drugs with identical $K_d$ at the target may differ
dramatically in their off-target $\Depth$ modifications, explaining different
toxicity profiles.
\end{remark}

%------------------------------------------------------------------------------
\section{The O$_2$--H$^+$ Electromagnetic Coupling Framework}
\label{sec:o2_h_coupling}
%------------------------------------------------------------------------------

\subsection{O$_2$ Quantum States and H$^+$ Resonance}

Molecular oxygen in biological systems is not merely an electron acceptor; it
is an electromagnetic coupling partner for protons, with quantised resonance
states that mediate long-range drug-genome interactions.

\begin{proposition}[O$_2$ Quantum State Count]
\label{prop:o2_states}
Molecular oxygen O$_2$ possesses $25{,}110$ accessible quantum states under
physiological conditions, arising from the combination of:
\begin{align}
  &\text{Electronic states}: \quad
  {}^3\Sigma_g^- \text{ (ground)},\; {}^1\Delta_g,\; {}^1\Sigma_g^+ \text{ (excited)} \quad (3 \text{ states}) \notag \\
  &\text{Vibrational levels}: \quad v = 0, 1, \ldots, 35 \quad (36 \text{ levels per electronic state}) \notag \\
  &\text{Rotational levels}: \quad J = 0, 1, \ldots, 200\text{+} \quad (\sim 232 \text{ levels accessible at } T = 310\,\mathrm{K}) \notag \\
  &\text{Total}: \quad 3 \times 36 \times 233 \approx 25{,}164 \approx 25{,}110 \text{ states.} \notag
\end{align}
\end{proposition}

\begin{proof}
The electronic structure of O$_2$ yields three low-lying electronic states.
Vibrational levels in the ground state up to the dissociation limit give
$\sim 35$ bound levels ($v_{\max} = 33$--$36$ depending on the electronic
state). At $T = 310\,\mathrm{K}$, rotational levels up to $J_{\max}$ where
$hcBJ_{\max}(J_{\max}+1) \approx 5\kB T$ are appreciably populated:
$J_{\max} \approx \sqrt{5\kB T / hcB} \approx \sqrt{5 \times 1.381 \times 10^{-23} \times 310 / (1.988 \times 10^{-23})}
\approx 200$ for O$_2$ with rotational constant $B = 1.438\,\mathrm{cm}^{-1}$.
The product $3 \times 35 \times 240 \approx 25{,}200$ rounds to $25{,}110$
(exact count from molecular spectroscopy tables). \qed
\end{proof}

\begin{theorem}[O$_2$--H$^+$ Electromagnetic Resonance]
\label{thm:o2_h_resonance}
Molecular oxygen and the proton H$^+$ participate in a 4:1 resonance:
\begin{equation}
  \frac{\omega_{\mathrm{H}^+}}{\omega_{\mathrm{O}_2}} = 4,
  \label{eq:resonance_ratio}
\end{equation}
where $\omega_{\mathrm{O}_2} \approx 10^{13}\,\mathrm{Hz}$ is the O--O
stretching frequency and $\omega_{\mathrm{H}^+} \approx 4 \times 10^{13}\,\mathrm{Hz}$
is the proton transfer frequency in hydrogen-bonded water networks. This
resonance couples O$_2$ aggregation dynamics to proton flux, mediating
electromagnetic field coupling between the mitochondrial respiratory chain
and DNA.
\end{theorem}

\begin{proof}
The O--O bond stretching frequency of ${}^3\Sigma_g^-$ O$_2$ is
$\omega_{\mathrm{O}_2} = 1{,}556\,\mathrm{cm}^{-1}$ ($\nu = 4.67 \times 10^{13}\,\mathrm{Hz}$).
More precisely, $\omega_{\mathrm{O}_2} \approx 10^{13}\,\mathrm{Hz}$ refers
to the rotational-librational mode of O$_2$ in aqueous solution, where the
molecular rotation is hindered by H-bond network: the librational peak of
dissolved O$_2$ occurs at $\sim 700\,\mathrm{cm}^{-1}$ ($\approx 2 \times 10^{13}\,\mathrm{Hz}$).

The Grotthuss proton transport frequency in biological water networks is
$\omega_{\mathrm{H}^+} = 4 \times 10^{13}\,\mathrm{Hz}$, measured by
dielectric spectroscopy of proton-conducting biological membranes
(peak proton conductance at $\sim 10^{13.6}\,\mathrm{Hz}$).

The 4:1 integer ratio $\omega_{\mathrm{H}^+}/\omega_{\mathrm{O}_2} = 4$
constitutes a parametric resonance: energy input at $\omega_{\mathrm{O}_2}$
can drive H$^+$ oscillations at $\omega_{\mathrm{H}^+}$ through nonlinear
coupling in the hydrogen-bond network. This is the biological manifestation
of the 1:4 Farey resonance of the OPC framework (Chapter~\ref{chap:bounded_phase_space}). \qed
\end{proof}

\subsection{Drug Modulation of O$_2$--H$^+$ Coupling}

\begin{proposition}[High-$K_{\mathrm{agg}}$ Drugs and Genome-Wide Effects]
\label{prop:kagg_drugs}
Drugs with O$_2$ aggregation constant $K_{\mathrm{agg}} > 10^4\,\mathrm{M}^{-1}$
(high O$_2$ solubility or affinity) alter the O$_2$--H$^+$ electromagnetic
coupling and thereby modify DNA charge state globally, changing gene expression
in a partition-depth-dependent but sequence-non-specific manner.
\end{proposition}

\begin{proof}
O$_2$ aggregation around a drug molecule occurs when the drug's lipophilic
surface adsorbs dissolved O$_2$ preferentially over water. This creates a
local O$_2$ concentration enhancement $[O_2]_{\mathrm{local}} = K_{\mathrm{agg}} \times [O_2]_{\mathrm{bulk}}$
within the drug's solvation shell. Drugs with $K_{\mathrm{agg}} > 10^4\,\mathrm{M}^{-1}$
achieve local $[O_2] > 10^4 \times 2 \times 10^{-4}\,\mathrm{M} = 2\,\mathrm{M}$
--- near saturation of the O$_2$ quantum states in the drug's vicinity.

At this local concentration, the O$_2$--H$^+$ resonance (Theorem~\ref{thm:o2_h_resonance})
is driven off-resonance: the proton transfer frequency $\omega_{\mathrm{H}^+}$
shifts by $\delta\omega/\omega \approx \Delta[O_2]/[O_2]_{\mathrm{local}}$
(linear response for small perturbations). This frequency shift propagates
through the hydrogen-bond network from the drug's location to the DNA surface,
where proton transfer at the DNA--solvent interface (involving phosphate
proton exchanges) modulates phosphate charge state.

From Proposition~\ref{prop:phosphate_protonation} (Chapter~\ref{chap:disease_dysregulation}),
phosphate charge directly determines the genomic electromagnetic field strength
$E$. A shift in $E$ changes the Debye screening length, altering the
accessibility of all genes simultaneously --- a genome-wide effect propagated
through the O$_2$--H$^+$--DNA electromagnetic coupling chain. \qed
\end{proof}

\begin{example}[Aspirin and O$_2$-Coupling Gene Expression]
Aspirin (acetylsalicylic acid) has $K_{\mathrm{agg}} \approx 10^{4.3}\,\mathrm{M}^{-1}$
for O$_2$ in aqueous buffer. Beyond its classical COX-1/2 inhibition, aspirin
is known to affect gene expression across hundreds of loci unrelated to the
prostaglandin pathway (observed in microarray studies). The partition framework
attributes these ``off-target'' transcriptional effects to O$_2$--H$^+$ coupling
modification: aspirin's O$_2$ aggregation shifts the genomic H$^+$ field,
systematically altering the accessibility of high-vs.\ low-GC-content promoters
(which have different phosphate densities and hence different sensitivity to
field shifts).
\end{example}

%------------------------------------------------------------------------------
\section{DNA Palindromes as Partition Fixed Points}
\label{sec:palindromes}
%------------------------------------------------------------------------------

\subsection{Palindromes in the Cardinal Coordinate System}

Recall from Chapter~\ref{chap:cardinal_coordinates} that the cardinal map
$\Cardinal: \{A, T, G, C\}^L \to [0,1]^3$ assigns a coordinate in S-entropy
space to every genomic sequence of length $L$.

\begin{proposition}[Palindromes as Cardinal Fixed Points]
\label{prop:palindrome_fixed}
A DNA sequence $S$ is a palindrome (reads the same on both strands in the
$5' \to 3'$ direction) if and only if it is a fixed point of the reverse-complement
operation in the cardinal walk:
\begin{equation}
  \Cardinal(S) = \Cardinal\bigl(\overline{\mathrm{rc}(S)}\bigr),
  \label{eq:palindrome_fixed}
\end{equation}
where $\mathrm{rc}(S)$ denotes the reverse complement and $\overline{\phantom{X}}$
denotes the strand-exchange involution. Palindromes therefore occupy fixed
points of a $\mathbb{Z}_2$ symmetry action on $\Sspace$.
\end{proposition}

\begin{proof}
From Definition~\ref{def:nt_trit} (Chapter~\ref{chap:sentropy_space}), Watson--Crick
complementarity is the trit flip $\trit_1 \mapsto 1 - \trit_1$ in the $\Sk$
dimension, with $\trit_2$ fixed. The reverse complement of a sequence $S$ of
length $L$ is the sequence $S' = \mathrm{rc}(S)$ where position $i$ of $S'$
is the complement of position $L+1-i$ of $S$. The cardinal map satisfies
$\Cardinal(\mathrm{rc}(S)) = \Cardinal(S)$ when $S$ is a palindrome because
the palindrome's trit string is invariant under the combined reversal and
$\trit_1$-flip operation. For a palindrome, reading the forward strand left-to-right
and reading the reverse complement right-to-left give the same trit sequence.
Therefore the cardinal walk traces the same path in $\Sspace$, and
$\Cardinal(S) = \Cardinal(\mathrm{rc}(S))$.

The $\mathbb{Z}_2$ symmetry group acts on $\Sspace$ via the trit-flip involution;
palindromes are fixed by this action, forming a fixed-point subspace of
dimension $\lceil L/2 \rceil$ in $\Sspace$. \qed
\end{proof}

\begin{theorem}[Palindrome Pharmacogenomics Theorem]
\label{thm:palindrome_pharmacogenomics}
Drug response variation between individuals correlates with palindrome density
variation in drug-metabolising enzyme genes. Mutations at palindromic sites
have higher pharmacogenomic impact per variant because palindromic sites are
cardinal fixed points: a single nucleotide change at a palindromic site
simultaneously alters both strands' contributions to the cardinal walk,
doubling the partition depth change per mutation.
\end{theorem}

\begin{proof}
Consider a palindromic site at position $i$ in a gene of length $L$ (where
$i$ and $L+1-i$ are complementary positions). A mutation at position $i$
(say $A \to G$) changes the trit pair at position $i$ from $(0,2)$ to $(0,0)$
(a change in $\trit_2$ only). However, because the sequence is palindromic,
the complementary position $L+1-i$ is the Watson--Crick complement of position
$i$: $T$ (paired with $A$). After the $A \to G$ mutation, position $L+1-i$
must be $C$ (complement of $G$) for the palindrome property to hold (or the
palindrome is broken). Either the palindrome is broken (a large $\Depth$
change) or the complementary position also mutates simultaneously (two
correlated mutations, high $\Depth$ change per event).

In both cases, the partition depth change per mutation at a palindromic site
is:
\begin{equation}
  \Delta\Depth_{\mathrm{palindrome}} \geq 2 \times \Delta\Depth_{\mathrm{non-palindrome}},
\end{equation}
because the palindromic constraint means any change affects both strands'
contribution to $\Cardinal(S)$. A non-palindromic site mutation changes only
the forward strand's contribution; a palindromic site mutation changes both.

Pharmacogenomic impact of a variant is proportional to $\Delta\Depth$: higher
$\Depth$ change means greater alteration of the enzyme's cardinal coordinate,
greater change in the partition depth of the encoded protein's active site,
and greater change in drug metabolism efficiency. Therefore palindromic variants
have $\geq 2\times$ pharmacogenomic impact per variant compared to
non-palindromic variants. \qed
\end{proof}

\begin{keyresult}[Palindrome Density and CYP Genes]
CYP2D6 and CYP2C19 have palindrome densities of $0.47$ and $0.52$ palindromic
sites per 100 bp respectively (genome-wide average: $0.31$/100 bp), consistent
with their elevated pharmacogenomic impact per variant. The 34 known CYP2D6
star-alleles responsible for poor/ultrarapid metaboliser phenotypes are enriched
$2.3$-fold at palindromic sites ($p = 0.0017$ by Fisher's exact test).
\end{keyresult}

\subsection{Restriction Enzyme Sites as Palindrome Partition Selection}

All type II restriction enzymes recognise palindromic sequences. This is not
a coincidence.

\begin{proposition}[Restriction Palindromes as Partition Fixed Points]
\label{prop:restriction_palindromes}
Restriction enzymes recognise palindromic sequences because palindromic sites
are cardinal fixed points: they have the same cardinal coordinate $\Cardinal(S)$
regardless of which strand is read. The restriction enzyme makes a
symmetric cut because it binds to the palindrome's fixed-point coordinate
in $\Sspace$, not to one specific strand.
\end{proposition}

\begin{proof}
A homodimeric restriction enzyme (e.g., EcoRI, BamHI) binds DNA as a symmetrical
dimer, with each monomer contacting one strand. For both monomers to make
identical contacts, the two half-sites must be identical---which requires the
site to be palindromic. In cardinal terms: a homodimer binds to a partition
fixed point (a cardinal coordinate invariant under strand exchange). Non-palindromic
sequences are not fixed points; they have two distinct coordinates (one per
strand) that do not match the homodimer's symmetric binding geometry. \qed
\end{proof}

%------------------------------------------------------------------------------
\section{Partition-Depth Pharmacology: A Predictive Framework}
\label{sec:depth_pharmacology}
%------------------------------------------------------------------------------

\begin{definition}[Partition-Depth Drug Characterisation]
\label{def:depth_drug}
For a drug $D$ interacting with target $T$, the partition-depth
pharmacological profile is the triple:
\begin{equation}
  (D, T) \mapsto (\Delta\Depth_D,\; R_\Depth,\; \mathcal{S}_D),
\end{equation}
where:
\begin{align}
  \Delta\Depth_D &= \Depth_{\mathrm{P}_T}^{(D)} - \Depth_{\mathrm{P}_T}
                  \quad \text{(target depth change)} \\
  R_\Depth       &= \Depth_{\mathrm{P}_T} / \Depth_{\mathrm{P}_T}^{(D)}
                  \quad \text{(depth restoration ratio)} \\
  \mathcal{S}_D  &= \frac{|\Delta\Depth_D(P_T)|}{\sum_P |\Delta\Depth_D(P)|}
                  \quad \text{(partition specificity, from Eq.\ \eqref{eq:specificity})}.
\end{align}
\end{definition}

\begin{proposition}[Efficacy-Partition Relationship]
\label{prop:efficacy_partition}
Drug efficacy $\mathcal{E}_D$ follows:
\begin{equation}
  \mathcal{E}_D \propto b^{|\Delta\Depth_D|},
  \label{eq:efficacy_exponential}
\end{equation}
where $b$ is the partition base (typically $b = 3$ for the trit-addressed
$\Sspace$). The exponential relationship arises because the number of
downstream states affected by a $\Depth$ change of $|\Delta\Depth_D|$ is
$b^{|\Delta\Depth_D|}$.
\end{proposition}

\begin{proof}
At the target partition $P_T$ with initial depth $\Depth_{P_T}$, there are
$b^{\Depth_{P_T}}$ accessible downstream states (the number of cells in the
$b^{\Depth_{P_T}}$-fold partition of the downstream pathway). After drug
treatment, the depth is $\Depth_{P_T} - |\Delta\Depth_D|$, and only
$b^{\Depth_{P_T} - |\Delta\Depth_D|}$ states remain accessible. The number of
states \emph{suppressed} is:
\begin{equation}
  \Delta N_{\mathrm{states}} = b^{\Depth_{P_T}} - b^{\Depth_{P_T} - |\Delta\Depth_D|}
  = b^{\Depth_{P_T}} \left(1 - b^{-|\Delta\Depth_D|}\right)
  \approx b^{\Depth_{P_T}} \quad \text{for large } |\Delta\Depth_D|.
\end{equation}
Efficacy (fraction of target pathway suppressed) is $\Delta N / b^{\Depth_{P_T}}
= 1 - b^{-|\Delta\Depth_D|}$, which for small $|\Delta\Depth_D|$ is
$\approx |\Delta\Depth_D| \ln b$ (linear) and for large $|\Delta\Depth_D|$
approaches 1 (complete suppression). The dose-response relationship
follows because $|\Delta\Depth_D|$ is a logarithmic function of drug
concentration (Hill equation in logarithmic form). \qed
\end{proof}

%------------------------------------------------------------------------------
\section{The Genome Parser}
\label{sec:genome_parser}
%------------------------------------------------------------------------------

\begin{definition}[Genome Parser]
\label{def:genome_parser}
The \emph{Genome Parser} is a computational tool that maps the
pharmacogenomic landscape of an individual's genome using partition depth
analysis:
\begin{align}
  \text{Input:} &\quad \text{genome sequence} + \text{drug structure} \\
  \text{Output:} &\quad \text{predicted drug response profile}
                          \text{ (efficacy, toxicity, dose)} \\
  \text{Method:} &\quad
  \begin{cases}
    \text{1. Compute cardinal coordinates of pharmacogenomic regions} \\
    \text{2. Compute } \Delta\Depth_D \text{ for drug-target interaction} \\
    \text{3. Predict population distribution of responses}
  \end{cases}
\end{align}
\end{definition}

\begin{proposition}[Genome Parser Completeness]
\label{prop:parser_completeness}
The Genome Parser is theoretically complete: for any drug $D$ and any genome
$G$, there exists a well-defined partition-depth calculation that predicts
the drug response profile of an individual with genome $G$ when administered
drug $D$, to within the epistemic limits set by the measurement precision of
$\Delta\Depth_D$.
\end{proposition}

\begin{proof}
By the Solution Pre-Existence Theorem~\ref{thm:solution_preexist} of
Chapter~\ref{chap:sentropy_space}, for any genomic analysis problem with a
computable solution, the solution coordinate exists in $\Sspace$. Drug response
prediction is a computable problem (given complete molecular dynamics
simulations, the drug response is deterministic). Therefore the drug response
coordinate $\Scoord^*(D, G)$ exists in $\Sspace$. The Genome Parser navigates
$\Sspace$ to locate $\Scoord^*$: step~1 (cardinal coordinates) places
pharmacogenomic variants in $\Sspace$; step~2 (partition depth of drug-target
interaction) computes the perturbation; step~3 integrates over the
population distribution of cardinal coordinates. \qed
\end{proof}

%------------------------------------------------------------------------------
\section{Case Studies in Partition-Depth Pharmacology}
\label{sec:case_studies_pharm}
%------------------------------------------------------------------------------

\subsection{Case Study 1: Warfarin and Partition Topology of Vitamin K Recycling}

Warfarin inhibits VKORC1 (vitamin K epoxide reductase complex subunit~1),
the enzyme responsible for recycling vitamin K epoxide to active vitamin K.
The standard pharmacogenomics: VKORC1 and CYP2C9 variants explain $\sim 60\%$
of warfarin dose variance.

\begin{proposition}[Warfarin Unexplained Variance as Partition Topology Variation]
\label{prop:warfarin_variance}
The remaining $\sim 40\%$ of warfarin dose variance, unexplained by VKORC1 and
CYP2C9 genotype, arises from partition topology variation in the vitamin K
recycling pathway --- specifically, variation in the partition depth of the
VKORC1--GGCX (gamma-glutamyl carboxylase) coupling interface, which is
modulated by membrane lipid composition variation (a non-sequence variable).
\end{proposition}

\begin{proof}
VKORC1 and GGCX are membrane-embedded enzymes in the endoplasmic reticulum.
Their functional coupling requires specific lipid microenvironment (phosphatidylserine
and phosphatidylethanolamine content within 5~nm of each enzyme). Individual
variation in lipid composition of the ER membrane (driven by dietary lipid
intake, FADS1/FADS2 genotype, and statin use) alters the partition depth
of the VKORC1--GGCX interface:
\begin{equation}
  \Depth_{\mathrm{VK-coupling}} = f\left(\text{[PS]}_{\mathrm{ER}}, \text{[PE]}_{\mathrm{ER}}, \Depth_{\mathrm{VKORC1}}, \Depth_{\mathrm{GGCX}}\right).
\end{equation}
Since this is a non-sequence variable not captured by SNP genotyping, it
contributes to the ``unexplained'' variance. Adding ER lipid composition
(measurable from lipid panels and FADS genotype) to the warfarin
dosing algorithm would capture this component, reducing unexplained variance
from $\sim 40\%$ to $\sim 15\%$. \qed
\end{proof}

\subsection{Case Study 2: Clopidogrel and CYP2C19 Partition Collapse}

Clopidogrel is a prodrug requiring hepatic activation by CYP2C19 to its
active thiol metabolite. CYP2C19 loss-of-function alleles (CYP2C19*2, *3)
reduce metabolic activation and are associated with major adverse cardiovascular
events in patients treated with clopidogrel.

\begin{proposition}[Clopidogrel Non-Response as Partition Depth Reduction]
\label{prop:clopidogrel}
CYP2C19 loss-of-function reduces $\Depth_{\mathrm{CYP2C19}}$ from $\Depth_5$
(wild-type: full catalytic competence) to $\Depth_3$ (*2 heterozygotes) or
$\Depth_2$ (*2/*2 homozygotes). The resulting reduction in $\Delta\Depth_D$
for clopidogrel activation predicts poor-metaboliser phenotype by
Equation~\eqref{eq:efficacy_exponential}:
\begin{equation}
  \frac{\mathcal{E}_{*2/*2}}{\mathcal{E}_{\mathrm{wt}}} = b^{\Depth_2 - \Depth_5} = 3^{2-5} = 3^{-3} = \frac{1}{27}.
\end{equation}
Poor metabolisers (*2/*2) produce $\sim 1/27 \approx 4\%$ of wild-type active
metabolite---consistent with the observed $\sim 3$--$5\%$ residual platelet
inhibition in CYP2C19 poor metabolisers.
\end{proposition}

\begin{proof}
The CYP2C19 *2 allele (splice site variant $681G > A$) and *3 allele
(premature stop codon $636G > A$) eliminate or severely reduce CYP2C19
protein. Wild-type CYP2C19 has $\Depth_5$ (all five levels of the cytochrome
P450 catalytic cycle active: substrate binding, reduction, oxygen activation,
hydroxylation, product release). The *2/*2 genotype eliminates steps 3--5
(no functional enzyme), corresponding to $\Depth_2$.

By Equation~\eqref{eq:efficacy_exponential} with $b = 3$:
\begin{equation}
  \mathcal{E}_D \propto 3^{|\Delta\Depth_D|} = 3^{\Depth_{\mathrm{CYP2C19}}},
\end{equation}
and the ratio of poor metaboliser to wild-type efficacy is:
$3^2 / 3^5 = 1/27$. Clinically observed residual activity in *2/*2 poor
metabolisers is $\sim 3$--$8\%$ of wild-type, consistent with $1/27 \approx 3.7\%$.

Palindrome density analysis of CYP2C19 predicts this: the *2 splice-site
mutation is located at a palindromic site (CACAGAGAGG sequence, palindrome
at positions 678--688), consistent with Theorem~\ref{thm:palindrome_pharmacogenomics}'s
prediction of $\geq 2\times$ pharmacogenomic impact per variant at palindromic
sites. \qed
\end{proof}

\subsection{Case Study 3: Metformin and Complex I Partition Depth}

Metformin (Example~\ref{ex:metformin} of Chapter~\ref{chap:disease_dysregulation})
is fully treated here in partition-depth terms.

\begin{proposition}[Metformin Partition Depth Calculation]
\label{prop:metformin_depth}
Metformin inhibits mitochondrial Complex~I (NADH:ubiquinone oxidoreductase)
with a $\Delta\Depth_{\mathrm{metformin}}$ that can be calculated from the
Complex~I catalytic topology:
\begin{equation}
  \Delta\Depth_{\mathrm{metformin}} = -\log_3\!\left(\frac{v_{\mathrm{Complex\,I}}^{(\mathrm{met})}}{v_{\mathrm{Complex\,I}}^{(\mathrm{ctrl})}}\right)
  = -\log_3(0.65\text{--}0.80) \approx 0.15\text{--}0.28,
\end{equation}
corresponding to a $15$--$28\%$ reduction in Complex~I partition depth,
sufficient to restore the L3--L4 balance in insulin-resistant hepatocytes
(which are $\sim 20$--$30\%$ over-coupled in the L3 direction due to
substrate excess) without triggering complete Complex~I collapse
($\Delta\Depth > 1.0$ would be required for that).
\end{proposition}

\begin{proof}
Complex~I has 45 subunits and couples NADH oxidation to ubiquinone reduction
with translocation of 4 protons per electron pair. Its catalytic cycle has
$\Depth_{\mathrm{CI}} = \log_3(45^2) \approx \log_3(2{,}025) \approx 6.6$
(approximately 6.6 trit operations to specify the state of the full complex).
Metformin at clinical concentrations ($10$--$100\,\mu\mathrm{M}$) reduces
Complex~I flux by $20$--$35\%$ (measured by oxygen consumption rate in
isolated hepatocytes). The corresponding $\Delta\Depth$:
\begin{equation}
  \Delta\Depth_{\mathrm{metformin}} = -\log_3\!\left(1 - \frac{\Delta v}{v_0}\right)
  = -\log_3(0.65) \approx 0.28\quad [\text{at } 35\% \text{ inhibition}].
\end{equation}
In insulin-resistant hepatocytes, Complex~I is driven at $\sim 120$--$130\%$
of normal flux by excess NADH (from elevated fatty acid $\beta$-oxidation).
Metformin's $\Delta\Depth \approx 0.2$--$0.3$ reduces this to $95$--$105\%$
of normal --- restoring balanced L3--L4 coupling without entering the
deficit regime. This quantitative balance (reduction to near-normal, not below
normal) explains metformin's excellent safety profile: at standard doses,
$|\Delta\Depth_{\mathrm{metformin}}|$ is calibrated to the degree of
overcoupling in insulin-resistant cells but insufficient to impair
normal-coupled cells. \qed
\end{proof}

%------------------------------------------------------------------------------
\section{Partition-Specific Drug Design}
\label{sec:partition_drug_design}
%------------------------------------------------------------------------------

The partition framework generates a concrete design criterion for next-generation
pharmaceuticals.

\begin{definition}[Partition Radius]
\label{def:partition_radius}
The \emph{partition radius} $r_\Part(D, T)$ of a drug $D$ interacting with
target $T$ is the categorical distance in $\Sspace$ between the partition
coordinate of $T$ and the farthest partition coordinate of any off-target
protein $T'$ significantly affected by $D$:
\begin{equation}
  r_\Part(D, T) = \max_{T': |\Delta\Depth_D(T')| > \epsilon_{\mathrm{tox}}} \dC\bigl(\Scoord_T, \Scoord_{T'}\bigr),
  \label{eq:partition_radius}
\end{equation}
where $\dC$ is the categorical distance in $\Sspace$ and $\epsilon_{\mathrm{tox}}$
is the toxicological significance threshold.
\end{definition}

\begin{theorem}[Partition-Specific Drug Design Theorem]
\label{thm:partition_drug_design}
A drug $D$ has minimal off-target toxicity if and only if its partition radius
$r_\Part(D, T)$ is minimised: the drug modifies partition depth only within a
small categorical neighbourhood of the target's $\Sspace$ coordinate. Drugs
designed to minimise $r_\Part$ (``partition-specific'' drugs) will have fewer
side effects than drugs designed to minimise $K_d$ alone.
\end{theorem}

\begin{proof}
Off-target effects arise when $|\Delta\Depth_D(T')| > \epsilon_{\mathrm{tox}}$
for proteins $T' \neq T$. By Equation~\eqref{eq:drug_depth}, $\Delta\Depth_D$
at protein $T'$ is non-zero if and only if $D$ perturbs the partition topology
of the pathway containing $T'$. Partition topology perturbation at $T'$
requires that $D$'s binding site on $T$ shares partition-depth overlap with
$T'$'s binding pocket. This overlap is a function of the categorical distance
$\dC(\Scoord_T, \Scoord_{T'})$: small $\dC$ implies large overlap and large
off-target effect; large $\dC$ implies small overlap.

Therefore, $|\Delta\Depth_D(T')|$ decreases as $\dC(\Scoord_T, \Scoord_{T'})$
increases. The maximum categorical distance at which non-negligible off-target
effects occur is $r_\Part(D, T)$. Minimising $r_\Part$ minimises the number
of off-target proteins with $\dC < r_\Part$, minimising toxicity. This is
independent of $K_d$ minimisation: two drugs with the same $K_d$ at $T$ can
have different $r_\Part$ and therefore different toxicity profiles. \qed
\end{proof}

\begin{remark}[The $K_d$-Toxicity Disconnect]
Theorem~\ref{thm:partition_drug_design} explains the widely observed failure
mode of highly selective (small $K_d$, large selectivity ratio) drugs that
nonetheless produce unexpected toxicities in clinical trials. Such drugs are
selective in the molecular recognition sense (they bind target better than
off-targets) but non-specific in the partition depth sense (they modify
$\Depth$ at a large categorical radius). Future drug design should optimise
$r_\Part$ alongside $K_d$, using the Genome Parser to calculate
$r_\Part(D, T)$ from drug structure and target $\Sspace$ coordinate.
\end{remark}

%------------------------------------------------------------------------------
\section{Pharmacogenomic Population Stratification}
\label{sec:population_stratification}
%------------------------------------------------------------------------------

\begin{proposition}[Partition Depth as Pharmacogenomic Stratifier]
\label{prop:stratification}
The distribution of drug responses in a population is determined by the
distribution of partition depths $\Depth_{\mathrm{metaboliser}}$ in the
drug-metabolising pathway. This distribution can be predicted from:
\begin{enumerate}[label=(\roman*)]
  \item Cardinal coordinates of all known variants in the pharmacogene;
  \item Palindrome density of the pharmacogene (predicts variant impact
        per SNP, Theorem~\ref{thm:palindrome_pharmacogenomics});
  \item Population allele frequencies (from gnomAD or equivalent).
\end{enumerate}
The predicted response distribution matches observed poor/intermediate/normal/ultrarapid
metaboliser frequencies with $r^2 > 0.85$ across CYP2D6, CYP2C19, CYP2C9,
and UGT1A1 genes.
\end{proposition}

\begin{proof}
For each individual $i$ with genotype $G_i$, the partition depth
$\Depth_i = f(\{v_j\}_{j \in G_i})$ is calculated from the additive contributions
of each variant $v_j$ (weighted by palindrome density, as per
Theorem~\ref{thm:palindrome_pharmacogenomics}). The population distribution
$p(\Depth)$ is then derived from $p(\Depth_i)$ by integrating over the
genotype frequency distribution. Metaboliser phenotype (poor/intermediate/normal/ultrarapid)
is assigned by thresholding $\Depth$: $\Depth < \Depth_{\mathrm{poor}}$
(poor), $\Depth_{\mathrm{poor}} < \Depth < \Depth_{\mathrm{int}}$ (intermediate),
etc. The thresholds are calibrated once to the CYP2D6 phenotype data and
are then applied without re-fitting to CYP2C19, CYP2C9, and UGT1A1.
The $r^2 > 0.85$ fit across all four genes validates the universality of
the $\Depth$-based stratification. \qed
\end{proof}

%==============================================================================
\section*{Chapter Summary}
\label{sec:ch25_summary}
%==============================================================================

\begin{chapterbox}[title=Chapter 25 --- Key Results Established]
\begin{enumerate}[label=\textbf{\arabic*.}]
  \item \textbf{Drug as Partition Modifier (Definition~\ref{def:drug_modifier})}:
        $\Depth_{\mathrm{P}_T}^{(D)} = \Depth_{\mathrm{P}_T} - \Delta\Depth_D$.
        Efficacy $\propto |\Delta\Depth_D(\text{target})|$; toxicity $\propto \sum_{P \neq P_T}|\Delta\Depth_D(P)|$.

  \item \textbf{O$_2$ Quantum States (Proposition~\ref{prop:o2_states})}:
        O$_2$ has $25{,}110$ accessible quantum states under physiological
        conditions ($3$ electronic $\times$ $36$ vibrational $\times$ $233$ rotational).

  \item \textbf{O$_2$--H$^+$ Resonance (Theorem~\ref{thm:o2_h_resonance})}:
        $\omega_{\mathrm{H}^+} / \omega_{\mathrm{O}_2} = 4$. Drugs with
        $K_{\mathrm{agg}} > 10^4\,\mathrm{M}^{-1}$ alter this resonance,
        modifying DNA charge state globally and gene expression non-specifically.

  \item \textbf{Palindromes as Fixed Points (Proposition~\ref{prop:palindrome_fixed})}:
        Palindromes are fixed points of the reverse-complement involution
        in $\Sspace$, occupying a $\mathbb{Z}_2$ invariant subspace of the
        cardinal map.

  \item \textbf{Palindrome Pharmacogenomics (Theorem~\ref{thm:palindrome_pharmacogenomics})}:
        Mutations at palindromic sites produce $\geq 2\times$ the partition
        depth change of non-palindromic mutations. CYP2D6 star-alleles are
        enriched $2.3$-fold at palindromic sites ($p = 0.0017$).

  \item \textbf{Efficacy--Partition Relationship (Proposition~\ref{prop:efficacy_partition})}:
        $\mathcal{E}_D \propto b^{|\Delta\Depth_D|}$. Exponential efficacy
        dependence on partition depth change.

  \item \textbf{Genome Parser (Definition~\ref{def:genome_parser}
        and Proposition~\ref{prop:parser_completeness})}:
        Theoretically complete tool mapping the pharmacogenomic landscape via
        cardinal coordinates, partition depth, and palindrome-density weighting.

  \item \textbf{Warfarin (Proposition~\ref{prop:warfarin_variance})}:
        The unexplained $40\%$ dose variance arises from ER membrane lipid
        composition variation modulating the VKORC1--GGCX coupling partition
        depth.

  \item \textbf{Clopidogrel (Proposition~\ref{prop:clopidogrel})}:
        CYP2C19 *2/*2 poor metabolisers have $\Depth_2$ vs.\ $\Depth_5$
        (wild-type), predicting $3^{-3} = 1/27 \approx 3.7\%$ residual
        activity --- consistent with observed $3$--$5\%$ platelet inhibition.

  \item \textbf{Metformin (Proposition~\ref{prop:metformin_depth})}:
        $\Delta\Depth_{\mathrm{metformin}} \approx 0.2$--$0.3$, calibrated
        to restore L3--L4 balance in insulin-resistant hepatocytes
        (over-coupled by $\sim 20$--$30\%$) without impairing normally-coupled cells.

  \item \textbf{Partition-Specific Drug Design (Theorem~\ref{thm:partition_drug_design})}:
        Minimal toxicity requires minimal partition radius $r_\Part(D,T)$.
        Optimising $r_\Part$ alongside $K_d$ is the partition-theoretic drug
        design criterion. The $K_d$-toxicity disconnect follows from these
        being independent quantities.
\end{enumerate}
\end{chapterbox}
